\Needspace{4\baselineskip}
\begin{Example}
\label{ex:differential-operator-boundary-domains}
We now discuss the differential operators considered in Section~\ref{subsec:boundary-value-problems-adjoints}. We let $\Omega$ be a domain with smooth boundary in $\mathbb{R}^m$, and let $p\in\mathbb{N}$ be given. We let $B_j(\bm{x},D)$, $j=1,\ldots,p$, be differential operators of order less than $2p$ which are well defined for $\bm{x}\in\partial\Omega$. We then define the domain
\begin{equation}\label{eq:differential-operator-boundary-domain}
\mathcal{D}_B(L):=\{u\in C_b^{2p}(\Omega)\mid B(\bm{x},D)u(\bm{x})=0,\bm{x}\in\partial\Omega,j=1,\ldots,p\}.
\end{equation}
We then define the operator $(\mathcal{D}_B(L),L)$ from $L^2(\Omega)$ to $L^2(\Omega)$ by
\begin{equation}\label{eq:general-differential-operator-action}
L(\bm{x},D)u(\bm{x}):=\sum_{|\alpha|\leq 2p}a_\alpha(\bm{x})D^\alpha u(\bm{x}).
\end{equation}

\end{Example}

\Needspace{4\baselineskip}
\begin{Example}
\label{ex:laplacian-dirichlet-neumann-domains}
In particular, we consider the \textbf{Laplacian}\[\Delta u:=\dfrac{\partial^2u}{\partial x_1^2}+\cdots+\dfrac{\partial^2u}{\partial x_m^2}\]again as an operator from $L^2(\Omega)$ to $L^2(\Omega)$. (For the moment, we let $\Omega$ be a bounded domain.) We consider two types of domains for the operator, the first corresponding to Dirichlet data:
\begin{equation}\label{eq:laplacian-classical-dirichlet-domain}
\widetilde{\mathcal{D}}_D(\Delta):=\{u\in C_b^2(\Omega)\mid u(\bm{x})=0\text{ for }\bm{x}\in\partial\Omega\};
\end{equation}
and the second corresponding to Neumann data:
\begin{equation}\label{eq:laplacian-classical-neumann-domain}
\widetilde{\mathcal{D}}_N(\Delta):=\left\{u\in C_b^2(\Omega)\mid\dfrac{\partial u}{\partial n}(\bm{x})=0\text{ for }\bm{x}\in\partial\Omega\right\}.
\end{equation}
For the moment we cannot say too much about the invertibility of the operator other than to observe that the operator $(\widetilde{\mathcal{D}}_N(\Delta),\Delta)$ has a nontrivial nullspace: the constant functions. We will examine the invertibility of these operators in Chapter~\ref{chap:linear-elliptic-equations}.
\end{Example}
